\documentclass[10pt,a4paper]{article}
\usepackage{amsmath,amssymb,amsthm,graphicx,verbatim,fontspec,hyperref}
\setmainfont{lmroman10-regular.otf}
% Retained preamble for the Stacks Project chapter snapshot in this directory.
% The chapter's own preamble.tex is a copy of the shared Stacks Project preamble; the block
% below is sliced MECHANICALLY, line for line, from it. Nothing was added, removed or
% redefined; the wrapper class, the xr-hyper/hyperref wiring and the externaldocument list
% are omitted because the assembler supplies the document class.
%
% lines taken from preamble.tex: 3, 5, 9, 11, 12, 156, 157, 158, ... (51 lines total)
%
% package coverage restored from the chapter's own preamble (2026-10-13): xy, xy-2cell

\usepackage[all]{xy}
\xyoption{2cell}
\UseAllTwocells

\newenvironment{reference}{\comment}{\endcomment}
\newenvironment{slogan}{\comment}{\endcomment}
\newenvironment{history}{\comment}{\endcomment}
\theoremstyle{plain}
\newtheorem{theorem}[subsection]{Theorem}
\newtheorem{proposition}[subsection]{Proposition}
\newtheorem{lemma}[subsection]{Lemma}
\theoremstyle{definition}
\newtheorem{definition}[subsection]{Definition}
\newtheorem{example}[subsection]{Example}
\newtheorem{exercise}[subsection]{Exercise}
\newtheorem{situation}[subsection]{Situation}
\theoremstyle{remark}
\newtheorem{remark}[subsection]{Remark}
\newtheorem{remarks}[subsection]{Remarks}
\numberwithin{equation}{subsection}
\def\lim{\mathop{\mathrm{lim}}\nolimits}
\def\colim{\mathop{\mathrm{colim}}\nolimits}
\def\Spec{\mathop{\mathrm{Spec}}}
\def\Hom{\mathop{\mathrm{Hom}}\nolimits}
\def\Ext{\mathop{\mathrm{Ext}}\nolimits}
\def\SheafHom{\mathop{\mathcal{H}\!\mathit{om}}\nolimits}
\def\SheafExt{\mathop{\mathcal{E}\!\mathit{xt}}\nolimits}
\def\Sch{\mathit{Sch}}
\def\Mor{\mathop{\mathrm{Mor}}\nolimits}
\def\Ob{\mathop{\mathrm{Ob}}\nolimits}
\def\Sh{\mathop{\mathit{Sh}}\nolimits}
\def\NL{\mathop{N\!L}\nolimits}
\def\CH{\mathop{\mathrm{CH}}\nolimits}
\def\proetale{{pro\text{-}\acute{e}tale}}
\def\etale{{\acute{e}tale}}
\def\QCoh{\mathit{QCoh}}
\def\Ker{\mathop{\mathrm{Ker}}}
\def\Im{\mathop{\mathrm{Im}}}
\def\Coker{\mathop{\mathrm{Coker}}}
\def\Coim{\mathop{\mathrm{Coim}}}
\DeclareMathSymbol{\boxtimes}{\mathbin}{AMSa}{"02}
\def\QCohstack{\mathcal{QC}\!\mathit{oh}}
\def\Cohstack{\mathcal{C}\!\mathit{oh}}
\def\Spacesstack{\mathcal{S}\!\mathit{paces}}
\def\Quotfunctor{\mathrm{Quot}}
\def\Hilbfunctor{\mathrm{Hilb}}
\def\Curvesstack{\mathcal{C}\!\mathit{urves}}
\def\Polarizedstack{\mathcal{P}\!\mathit{olarized}}
\def\Complexesstack{\mathcal{C}\!\mathit{omplexes}}
\def\Pic{\mathop{\mathrm{Pic}}\nolimits}
\def\Picardstack{\mathcal{P}\!\mathit{ic}}
\def\Picardfunctor{\mathrm{Pic}}
\def\Deformationcategory{\mathcal{D}\!\mathit{ef}}
\title{Grothendieck spectral sequence for compositions of right derived functors}
\author{The Stacks Project Authors}
\date{stacks-project:derived.tex@a04446e57ec1fbc252a871afcec7752fb2807b14; GNU FDL 1.2 or later}
\begin{document}
\maketitle
\noindent\textit{Editorial note.} Original statement, complete author proof, Cartan–Eilenberg construction and both source-local spectral-sequence calculations retained. External original labels remain explicit source pointers, without guessed theorem numbering. Shared pinned my.bib supplied; no collector mathematical repairs.\par\medskip
\noindent\textit{Allowed standard prerequisites (editorial).}\begin{itemize}
\item Abelian categories with enough injectives, bounded-below derived functors and double-complex spectral-sequence notation. Complete source-local Cartan–Eilenberg construction, page computations and composition argument retained. Ordinary injective-resolution and first-quadrant convergence facts remain exactly cited standard prerequisites.
\end{itemize}\medskip
\section{Right derived functors and injective resolutions}
\label{section-right-derived-functor}

\noindent
At this point we can use the material above to define the right derived
functors of an additive functor between an abelian category having
enough injectives and a general abelian category.

\begin{lemma}
\label{lemma-injective-acyclic}
Let $\mathcal{A}$ be an abelian category.
Let $I \in \Ob(\mathcal{A})$ be an injective object.
Let $I^\bullet$ be a bounded below complex of injectives in $\mathcal{A}$.
\begin{enumerate}
\item $I^\bullet$ computes $RF$ relative to $\text{Qis}^{+}(\mathcal{A})$
for any exact functor $F : K^{+}(\mathcal{A}) \to \mathcal{D}$
into any triangulated category $\mathcal{D}$.
\item $I$ is right acyclic for any additive functor
$F : \mathcal{A} \to \mathcal{B}$ into any abelian category $\mathcal{B}$.
\end{enumerate}
\end{lemma}

\begin{proof}
Part (2) is a direct consequences of part (1) and
Definition [original label: \nolinkurl{definition-derived-functor}].
To prove (1) let $\alpha : I^\bullet \to K^\bullet$ be a quasi-isomorphism
into a complex. By
Lemma [original label: \nolinkurl{lemma-morphisms-lift}]
we see that $\alpha$ has a left inverse. Hence the category
$I^\bullet/\text{Qis}^{+}(\mathcal{A})$ is essentially constant with value
$\text{id} : I^\bullet \to I^\bullet$. Thus also the ind-object
$$
I^\bullet/\text{Qis}^{+}(\mathcal{A}) \longrightarrow \mathcal{D}, \quad
(I^\bullet \to K^\bullet) \longmapsto F(K^\bullet)
$$
is essentially constant with value $F(I^\bullet)$. This proves (1), see
Definitions [original label: \nolinkurl{definition-right-derived-functor-defined}] and
[original label: \nolinkurl{definition-computes}].
\end{proof}

\begin{lemma}
\label{lemma-enough-injectives-right-derived}
Let $\mathcal{A}$ be an abelian category with enough injectives.
\begin{enumerate}
\item For any exact functor $F : K^{+}(\mathcal{A}) \to \mathcal{D}$
into a triangulated category $\mathcal{D}$ the right derived
functor
$$
RF : D^{+}(\mathcal{A}) \longrightarrow \mathcal{D}
$$
is everywhere defined.
\item For any additive functor $F : \mathcal{A} \to \mathcal{B}$ into an
abelian category $\mathcal{B}$ the right derived functor
$$
RF : D^{+}(\mathcal{A}) \longrightarrow D^{+}(\mathcal{B})
$$
is everywhere defined.
\end{enumerate}
\end{lemma}

\begin{proof}
Combine
Lemma \ref{lemma-injective-acyclic}
and
Proposition [original label: \nolinkurl{proposition-enough-acyclics}]
for the second assertion. To see the first assertion combine
Lemma [original label: \nolinkurl{lemma-injective-resolutions-exist}],
Lemma \ref{lemma-injective-acyclic},
Lemma [original label: \nolinkurl{lemma-existence-computes}],
and Equation ([original label: \nolinkurl{equation-everywhere}]).
\end{proof}

\begin{lemma}
\label{lemma-right-derived-properties}
Let $\mathcal{A}$ be an abelian category with enough injectives.
Let $F : \mathcal{A} \to \mathcal{B}$ be an additive functor.
\begin{enumerate}
\item The functor $RF$ is an exact functor
$D^{+}(\mathcal{A}) \to D^{+}(\mathcal{B})$.
\item The functor $RF$ induces an exact functor
$K^{+}(\mathcal{A}) \to D^{+}(\mathcal{B})$.
\item The functor $RF$ induces a $\delta$-functor
$\text{Comp}^{+}(\mathcal{A}) \to D^{+}(\mathcal{B})$.
\item The functor $RF$ induces a $\delta$-functor
$\mathcal{A} \to D^{+}(\mathcal{B})$.
\end{enumerate}
\end{lemma}

\begin{proof}
This lemma simply reviews some of the results obtained so far.
Note that by
Lemma \ref{lemma-enough-injectives-right-derived}
$RF$ is everywhere defined. Here are some references:
\begin{enumerate}
\item The derived functor is exact: This boils down to
Lemma [original label: \nolinkurl{lemma-2-out-of-3-defined}].
\item This is true because $K^{+}(\mathcal{A}) \to D^{+}(\mathcal{A})$
is exact and compositions of exact functors are exact.
\item This is true because
$\text{Comp}^{+}(\mathcal{A}) \to D^{+}(\mathcal{A})$ is
a $\delta$-functor, see
Lemma [original label: \nolinkurl{lemma-derived-canonical-delta-functor}].
\item This is true because $\mathcal{A} \to \text{Comp}^{+}(\mathcal{A})$
is exact and precomposing a $\delta$-functor by an exact functor gives
a $\delta$-functor.
\end{enumerate}
\end{proof}

\begin{lemma}
\label{lemma-higher-derived-functors}
Let $\mathcal{A}$ be an abelian category with enough injectives.
Let $F : \mathcal{A} \to \mathcal{B}$ be a left exact functor.
\begin{enumerate}
\item For any short exact sequence
$0 \to A^\bullet \to B^\bullet \to C^\bullet \to 0$
of complexes in $\text{Comp}^{+}(\mathcal{A})$ there
is an associated long exact sequence
$$
\ldots \to
H^i(RF(A^\bullet)) \to
H^i(RF(B^\bullet)) \to
H^i(RF(C^\bullet)) \to
H^{i + 1}(RF(A^\bullet)) \to \ldots
$$
\item The functors $R^iF : \mathcal{A} \to \mathcal{B}$
are zero for $i < 0$. Also $R^0F = F : \mathcal{A} \to \mathcal{B}$.
\item We have $R^iF(I) = 0$ for $i > 0$ and $I$ injective.
\item The sequence $(R^iF, \delta)$ forms a universal $\delta$-functor (see
Homology, Definition [original label: \nolinkurl{homology-definition-universal-delta-functor}])
from $\mathcal{A}$ to $\mathcal{B}$.
\end{enumerate}
\end{lemma}

\begin{proof}
This lemma simply reviews some of the results obtained so far.
Note that by
Lemma \ref{lemma-enough-injectives-right-derived}
$RF$ is everywhere defined. Here are some references:
\begin{enumerate}
\item This follows from
Lemma \ref{lemma-right-derived-properties} part (3)
combined with the long exact cohomology sequence
([original label: \nolinkurl{equation-long-exact-cohomology-sequence-D}]) for
$D^{+}(\mathcal{B})$.
\item This is
Lemma [original label: \nolinkurl{lemma-left-exact-higher-derived}].
\item This is the fact that injective objects are acyclic.
\item This is
Lemma [original label: \nolinkurl{lemma-right-derived-delta-functor}].
\end{enumerate}
\end{proof}







\section{Cartan-Eilenberg resolutions}
\label{section-cartan-eilenberg}

\noindent
This section can be expanded. The material can be generalized and applied in
more cases. Resolutions need not use injectives and the method also
works in the unbounded case in some situations.

\begin{definition}
\label{definition-cartan-eilenberg}
Let $\mathcal{A}$ be an abelian category.
Let $K^\bullet$ be a bounded below complex.
A {\it Cartan-Eilenberg resolution} of $K^\bullet$
is given by a double complex $I^{\bullet, \bullet}$
and a morphism of complexes $\epsilon : K^\bullet \to I^{\bullet, 0}$
with the following properties:
\begin{enumerate}
\item There exists a $i \ll 0$ such that $I^{p, q} = 0$ for all $p < i$
and all $q$.
\item We have $I^{p, q} = 0$ if $q < 0$.
\item The complex $I^{p, \bullet}$ is an injective resolution of $K^p$.
\item The complex $\Ker(d_1^{p, \bullet})$ is an injective resolution
of $\Ker(d_K^p)$.
\item The complex $\Im(d_1^{p, \bullet})$ is an injective resolution
of $\Im(d_K^p)$.
\item The complex $H^p_I(I^{\bullet, \bullet})$ is an injective resolution
of $H^p(K^\bullet)$.
\end{enumerate}
\end{definition}

\begin{lemma}
\label{lemma-cartan-eilenberg}
Let $\mathcal{A}$ be an abelian category with enough injectives.
Let $K^\bullet$ be a bounded below complex.
There exists a Cartan-Eilenberg resolution of $K^\bullet$.
\end{lemma}

\begin{proof}
Suppose that $K^p = 0$ for $p < n$. Decompose $K^\bullet$ into
short exact sequences as follows: Set $Z^p = \Ker(d^p)$,
$B^p = \Im(d^{p - 1})$, $H^p = Z^p/B^p$, and consider
$$
\begin{matrix}
0 \to Z^n \to K^n \to B^{n + 1} \to 0 \\
0 \to B^{n + 1} \to Z^{n + 1} \to H^{n + 1} \to 0 \\
0 \to Z^{n + 1} \to K^{n + 1} \to B^{n + 2} \to 0 \\
0 \to B^{n + 2} \to Z^{n + 2} \to H^{n + 2} \to 0 \\
\ldots
\end{matrix}
$$
Set $I^{p, q} = 0$ for $p < n$. Inductively we choose
injective resolutions as follows:
\begin{enumerate}
\item Choose an injective resolution $Z^n \to J_Z^{n, \bullet}$
with $J_Z^{n, m} = 0$ for $m < 0$, see
Lemma [original label: \nolinkurl{lemma-injective-resolutions-exist}].
\item Using Lemma [original label: \nolinkurl{lemma-injective-resolution-ses}] choose injective
resolutions $K^n \to I^{n, \bullet}$, $B^{n + 1} \to J_B^{n + 1, \bullet}$
with $I^{n, m} = 0$ for $m < 0$ and $J_B^{n + 1, m} = 0$ for $m < 0$,
and an exact sequence of complexes
$0 \to J_Z^{n, \bullet} \to I^{n, \bullet} \to J_B^{n + 1, \bullet} \to 0$
compatible with the short exact sequence
$0 \to Z^n \to K^n \to B^{n + 1} \to 0$.
\item Using Lemma [original label: \nolinkurl{lemma-injective-resolution-ses}] choose injective
resolutions $Z^{n + 1} \to J_Z^{n + 1, \bullet}$,
$H^{n + 1} \to J_H^{n + 1, \bullet}$
with $J_Z^{n + 1, m} = 0$ for $m < 0$ and $J_H^{n + 1, m} = 0$ for $m < 0$
and an exact sequence of complexes
$0 \to J_B^{n + 1, \bullet} \to J_Z^{n + 1, \bullet}
\to J_H^{n + 1, \bullet} \to 0$
compatible with the short exact sequence
$0 \to B^{n + 1} \to Z^{n + 1} \to H^{n + 1} \to 0$.
\item Etc.
\end{enumerate}
Taking as maps $d_1^\bullet : I^{p, \bullet} \to I^{p + 1, \bullet}$
the compositions
$I^{p, \bullet} \to J_B^{p + 1, \bullet} \to
J_Z^{p + 1, \bullet} \to I^{p + 1, \bullet}$ everything is clear.
\end{proof}

\begin{lemma}
\label{lemma-two-ss-complex-functor}
Let $F : \mathcal{A} \to \mathcal{B}$ be a left exact functor of
abelian categories.
Let $K^\bullet$ be a bounded below complex of $\mathcal{A}$.
Let $I^{\bullet, \bullet}$ be a Cartan-Eilenberg resolution
for $K^\bullet$. The spectral sequences
$({}'E_r, {}'d_r)_{r \geq 0}$ and $({}''E_r, {}''d_r)_{r \geq 0}$
associated to the double complex $F(I^{\bullet, \bullet})$
satisfy the relations
$$
{}'E_1^{p, q} = R^qF(K^p)
\quad
\text{and}
\quad
{}''E_2^{p, q} = R^pF(H^q(K^\bullet))
$$
Moreover, these spectral sequences are bounded, converge to
$H^*(RF(K^\bullet))$, and the associated induced filtrations on
$H^n(RF(K^\bullet))$ are finite.
\end{lemma}

\begin{proof}
We will use the following remarks without further mention:
\begin{enumerate}
\item As $I^{p, \bullet}$ is an injective resolution of
$K^p$ we see that $RF$ is defined at $K^p[0]$
with value $F(I^{p, \bullet})$.
\item As $H^p_I(I^{\bullet, \bullet})$ is an injective resolution of
$H^p(K^\bullet)$ the derived functor $RF$ is defined at $H^p(K^\bullet)[0]$
with value $F(H^p_I(I^{\bullet, \bullet}))$.
\item By
Homology, Lemma [original label: \nolinkurl{homology-lemma-double-complex-gives-resolution}]
the total complex $\text{Tot}(I^{\bullet, \bullet})$
is an injective resolution of
$K^\bullet$. Hence $RF$ is defined at $K^\bullet$ with value
$F(\text{Tot}(I^{\bullet, \bullet}))$.
\end{enumerate}
Consider the two spectral sequences associated to the double complex
$L^{\bullet, \bullet} = F(I^{\bullet, \bullet})$, see
Homology, Lemma [original label: \nolinkurl{homology-lemma-ss-double-complex}].
These are both bounded, converge to $H^*(\text{Tot}(L^{\bullet, \bullet}))$,
and induce finite filtrations on $H^n(\text{Tot}(L^{\bullet, \bullet}))$, see
Homology, Lemma [original label: \nolinkurl{homology-lemma-first-quadrant-ss}].
Since
$\text{Tot}(L^{\bullet, \bullet}) =
\text{Tot}(F(I^{\bullet, \bullet})) =
F(\text{Tot}(I^{\bullet, \bullet}))$ computes
$RF(K^\bullet)$ we find the final assertion of the lemma holds true.

\medskip\noindent
Computation of the first spectral sequence. We have
${}'E_1^{p, q} = H^q(L^{p, \bullet})$ in other words
$$
{}'E_1^{p, q} = H^q(F(I^{p, \bullet})) = R^qF(K^p)
$$
as desired. Observe for later use that the maps
${}'d_1^{p, q} : {}'E_1^{p, q} \to {}'E_1^{p + 1, q}$ are the maps
$R^qF(K^p) \to R^qF(K^{p + 1})$ induced by $K^p \to K^{p + 1}$
and the fact that $R^qF$ is a functor.

\medskip\noindent
Computation of the second spectral sequence. We have
${}''E_1^{p, q} = H^q(L^{\bullet, p}) = H^q(F(I^{\bullet, p}))$.
Note that the complex $I^{\bullet, p}$ is bounded below,
consists of injectives, and moreover each kernel, image, and
cohomology group of the differentials is an injective object
of $\mathcal{A}$. Hence we can split the differentials, i.e.,
each differential is a split surjection onto a direct summand.
It follows that the same is true after applying $F$. Hence
${}''E_1^{p, q} = F(H^q(I^{\bullet, p})) = F(H^q_I(I^{\bullet, p}))$.
The differentials on this are $(-1)^q$ times $F$ applied to
the differential of the complex $H^p_I(I^{\bullet, \bullet})$
which is an injective resolution of $H^p(K^\bullet)$. Hence the
description of the $E_2$ terms.
\end{proof}

\begin{remark}
\label{remark-functorial-ss}
The spectral sequences of Lemma \ref{lemma-two-ss-complex-functor}
are functorial in the complex $K^\bullet$. This follows from functoriality
properties of Cartan-Eilenberg resolutions. On the other hand, they are
both examples of a more general spectral sequence which may be associated
to a filtered complex of $\mathcal{A}$. The functoriality will follow from
its construction. We will return to this in the section on the filtered
derived category, see Remark [original label: \nolinkurl{remark-final-functorial}].
\end{remark}










\section{Composition of right derived functors}

\label{section-composition-right-derived-functors}

\noindent
Sometimes we can compute the right derived functor of a composition.
Suppose that $\mathcal{A}, \mathcal{B}, \mathcal{C}$ be abelian categories.
Let $F : \mathcal{A} \to \mathcal{B}$ and $G : \mathcal{B} \to \mathcal{C}$
be left exact functors. Assume that the right derived functors
$RF : D^{+}(\mathcal{A}) \to D^{+}(\mathcal{B})$,
$RG : D^{+}(\mathcal{B}) \to D^{+}(\mathcal{C})$, and
$R(G \circ F) : D^{+}(\mathcal{A}) \to D^{+}(\mathcal{C})$
are everywhere defined. Then there exists a canonical transformation
$$
t : R(G \circ F) \longrightarrow RG \circ RF
$$
of functors from $D^{+}(\mathcal{A})$ to $D^{+}(\mathcal{C})$, see
Lemma [original label: \nolinkurl{lemma-compose-derived-functors-general}].
This transformation need not always be an isomorphism.

\begin{lemma}
\label{lemma-compose-derived-functors}
Let $\mathcal{A}, \mathcal{B}, \mathcal{C}$ be abelian categories.
Let $F : \mathcal{A} \to \mathcal{B}$ and $G : \mathcal{B} \to \mathcal{C}$
be left exact functors. Assume $\mathcal{A}$, $\mathcal{B}$ have
enough injectives. The following are equivalent
\begin{enumerate}
\item $F(I)$ is right acyclic for $G$ for each injective object $I$
of $\mathcal{A}$, and
\item the canonical map
$$
t : R(G \circ F) \longrightarrow RG \circ RF.
$$
is isomorphism of functors from $D^{+}(\mathcal{A})$ to $D^{+}(\mathcal{C})$.
\end{enumerate}
\end{lemma}

\begin{proof}
If (2) holds, then (1) follows by evaluating the isomorphism
$t$ on $RF(I) = F(I)$. Conversely, assume (1) holds.
Let $A^\bullet$ be a bounded below complex of $\mathcal{A}$.
Choose an injective resolution $A^\bullet \to I^\bullet$.
The map $t$ is given (see proof of
Lemma [original label: \nolinkurl{lemma-compose-derived-functors-general}])
by the maps
$$
R(G \circ F)(A^\bullet) =
(G \circ F)(I^\bullet) =
G(F(I^\bullet))) \to
RG(F(I^\bullet)) =
RG(RF(A^\bullet))
$$
where the arrow is an isomorphism by
Lemma [original label: \nolinkurl{lemma-leray-acyclicity}].
\end{proof}


\begin{lemma}[Grothendieck spectral sequence]
\label{lemma-grothendieck-spectral-sequence}
With assumptions as in Lemma \ref{lemma-compose-derived-functors}
and assuming the equivalent conditions (1) and (2) hold.
Let $X$ be an object of $D^{+}(\mathcal{A})$.
There exists a spectral sequence $(E_r, d_r)_{r \geq 0}$
consisting of bigraded objects $E_r$ of $\mathcal{C}$ with
$d_r$ of bidegree $(r, - r + 1)$ and with
$$
E_2^{p, q} = R^pG(H^q(RF(X)))
$$
Moreover, this spectral sequence is bounded, converges to
$H^*(R(G \circ F)(X))$, and induces a finite filtration
on each $H^n(R(G \circ F)(X))$.
\end{lemma}


\noindent
For an object $A$ of $\mathcal{A}$ we get
$E_2^{p, q} = R^pG(R^qF(A))$ converging to $R^{p + q}(G \circ F)(A)$.


\begin{proof}
We may represent $X$ by a bounded below complex $A^\bullet$.
Choose an injective resolution $A^\bullet \to I^\bullet$.
Choose a Cartan-Eilenberg resolution
$F(I^\bullet) \to I^{\bullet, \bullet}$ using
Lemma \ref{lemma-cartan-eilenberg}.
Apply the second spectral sequence of
Lemma \ref{lemma-two-ss-complex-functor}.
\end{proof}








\bibliographystyle{plain}
\bibliography{my}
\end{document}
